% TOP_PROBLEM: {"id": 30003263, "title": "Macroscopic Limits and Phase Diagrams of Ferroelectric Models", "category_id": 16, "category": {"id": 16, "name": "physics", "display_name": "Mathematical Physics", "description": "Problems at the intersection of mathematics and physics.", "slug": "physics", "order_index": 16, "created_at": "2026-07-31T15:26:25.671Z"}, "status": "open", "statusQualification": "The still-open lower-bound direction is expressly part of the first question, and the lead titanate--lead zirconate phase diagram is a second target. The microscopic and effective Hamiltonians are source-defined and are not replaced by a generic Ising model."}
% FIRST_POSTED: 2026-10-11
% ENTRY_KIND: statement_only
% SUBMISSION: {"kind": "statement_only", "category": "top_problems", "problem_id": "30003263", "proposal_type": "established", "source": "Mini-Workshop: Mathematics of Magnetoelastic Materials (2016), Oberwolfach Reports, DOI 10.4171/owr/2016/51, pp. 2909--2940. Primary report accessed 2026-10-11; relevant PDF pages 25. https://doi.org/10.4171/owr/2016/51; EMS report text (catalogue-listed context; not independently consulted in this preparation). https://ems.press/content/serial-article-files/46654?nt=1", "status_evidence": "Canonical source status open retained. Supplied literature triage dated 2026-08-21: The cited work supplies an upper bound and numerical evidence, but explicitly leaves the lower bound and systematic analysis of the proposed long-range random-field model unresolved. Primary-report access checked 2026-10-11; this historical report and inherited triage do not certify current resolution of every component.", "status_checked_on": "2026-10-11", "duplicate_check": "Compared canonical ID, title and normalized target with the complete local UnsolvedMath registry; normalized exact-text matches: none; title matches: none. Related or overlapping canonical records: none identified in this batch. Primary source context and the existing OWR statement audit used for scope and duplicate review; an exhaustive mathematical-equivalence/literature search was not performed.", "attribution": "Original proposer not named in the supplied canonical record; attribution remains to the identified Oberwolfach contribution/report, not its editors as invented proposers. Explanatory statement preparation and source/semantic audit assisted by OpenAI GPT-6 Codex; no human mathematical authorship or independent proof verification is claimed."}
% SOURCE_STATEMENT: "Do macroscopic polarization models for ferroelectric solids arise as the large-body limit of microscopic electron-density models, including the still-open lower-bound direction? Can one also systematically analyze the proposed random-field-Ising-like model with long-range interactions for the phase diagram of lead titanate–lead zirconate solutions?"
% Preparation and final self-audit: CONTRIBUTING.md prompts C and D; source/audit artifacts retained in the local physics_statements_2026_10_11 batch.
% Imported from: physics_statements.tex; UnsolvedMath ID 30003263
% !TeX program = lualatex
\documentclass[11pt]{article}
\usepackage[margin=25mm]{geometry}
\usepackage{fontspec}
\setmainfont{Latin Modern Roman}
\usepackage{amsmath,amssymb,mathtools}
\usepackage{unicode-math}
\setmathfont{Latin Modern Math}
\usepackage{xurl,hyperref}
\hypersetup{colorlinks=true,urlcolor=blue}
\setcounter{secnumdepth}{0}
\setlength{\parindent}{0pt}
\setlength{\parskip}{5pt}
\setlength{\emergencystretch}{3em}
\newcommand{\sref}[2]{\hyperlink{source:#1:#2}{[S#2]}}
\newcommand{\cataloguescope}[1]{\paragraph{Recorded target.}#1}
\newcommand{\R}{\mathbb R}
\newcommand{\C}{\mathbb C}
\newcommand{\N}{\mathbb N}
\newcommand{\Z}{\mathbb Z}
\newcommand{\Q}{\mathbb Q}
\newcommand{\F}{\mathbb F}
\newcommand{\eps}{\varepsilon}
\begin{document}
% UnsolvedMath ID 30003263; source catalogue code OWR-15172-002
\setcounter{section}{30003262}
\section{Macroscopic Limits and Phase Diagrams of Ferroelectric Models}\label{30003263}
{\small UnsolvedMath ID 30003263 · Mathematical Physics}
\subsection{Definitions and mathematical statement}

\paragraph{Shared notation.}$\mathbb N=\{1,2,\ldots\}$; the standard number fields have their usual meanings. All additional conventions are stated locally.


\paragraph{Definitions and conventions.}
Polarization is the macroscopic dipole moment per unit volume. A large-body limit relates microscopic electron-density energies to a continuum polarization functional as system size grows. In a variational limit, a lower bound controls all convergent sequences, whereas an upper bound is produced by recovery configurations. A random-field Ising-like model has local polarization variables, random fields and pair interactions; long-range interactions remain significant over many lattice spacings.

\paragraph{Target retained from the canonical source catalogue.}
\paragraph{Catalogue wording retained for traceability.}
Do macroscopic polarization models for ferroelectric solids arise as the large-body limit of microscopic electron-density models, including the still-open lower-bound direction? Can one also systematically analyze the proposed random-field-Ising-like model with long-range interactions for the phase diagram of lead titanate–lead zirconate solutions?

\paragraph{Source scope and qualification.}
The still-open lower-bound direction is expressly part of the first question, and the lead titanate--lead zirconate phase diagram is a second target. The microscopic and effective Hamiltonians are source-defined and are not replaced by a generic Ising model.
\subsection{Short English statement}
Derive ferroelectric polarization models from microscopic theories and analyze the proposed long-range random-field phase diagram.
\subsection{Sources}
\begin{itemize}
\item[\hypertarget{source:30003263:1}{S1}] ULAM.AI and the UnsolvedMath Contributors, UnsolvedMath, record 30003263 (OWR-15172-002). Catalogue curation under CC BY 4.0; definitions and English exposition added locally.
\url{https://huggingface.co/datasets/ulamai/UnsolvedMath}.
\item[\hypertarget{source:30003263:2}{S2}] Mini-Workshop: Mathematics of Magnetoelastic Materials (2016), Oberwolfach Reports, DOI 10.4171/owr/2016/51, pp. 2909--2940. Primary report accessed 2026-10-11; relevant PDF pages 25.
\url{https://doi.org/10.4171/owr/2016/51}.
\item[\hypertarget{source:30003263:3}{S3}] EMS report text (catalogue-listed context; not independently consulted in this preparation).
\url{https://ems.press/content/serial-article-files/46654?nt=1}.
\end{itemize}
\label{end:30003263}
\end{document}
